\section{Keeping norm and content in the same Euler product}
\label{jct:zero-free-proof}
This section proves the general tail theorem.  Its extra input is
exactly the two zero-free half-planes stated in Section~1. The
specialization \(\beta=\rho=1\) requires no additional theorem.

\paragraph{The two complex variables.}
For \(\Re s>1\) and \(\Re v>1\), define the absolutely convergent series
\begin{equation}
 \mathcal F(s,v)=\sum_{\mathfrak a}
 m(\mathfrak a)(\N\mathfrak a)^{-s}
 \cont(\mathfrak a)^{2s-v}.
 \label{jct:bivariate}
\end{equation}
The powers of positive integers use their real logarithms.  Let
\(\eta=\chi_{-3}\), and at a rational prime \(p\ne3,269\) put
\[
 t_p=\begin{cases}\psi(\mathfrak p)+\psi(\overline{\mathfrak p}),
                   &p\text{ split in }K,\\0,&p\text{ inert},\end{cases}
 \qquad D_p(s)=1-t_pp^{-s}+\eta(p)p^{-2s}.
\]
Thus \(P(s)^{-1}=\prod_{p\ne3,269}D_p(s)\).
At a split prime the four squarefree ideal possibilities are
\(1,\mathfrak p,\overline{\mathfrak p},(p)\).
Their contributions to \eqref{jct:bivariate} are
\(1,-\psi(\mathfrak p)p^{-s},
-\psi(\overline{\mathfrak p})p^{-s},p^{-v}\).
At an inert prime the two possibilities contribute \(1,-p^{-v}\).
Consequently the local factor is, in both cases,
\begin{equation}
 1-t_pp^{-s}+\eta(p)p^{-v}.
 \label{jct:bivariate-local}
\end{equation}
This computation includes ideals of content divisible by primes
in the other ideal; no cross-coprimality is introduced.

Put
\begin{equation}
 A(v)=A_1(v)=\frac{L(v,\eta)}
 {\zeta(2v)(1-3^{-2v})(1-269^{-v})},\qquad
 E(s,v)=\prod_{p\ne3,269}
 \frac{1-t_pp^{-s}+\eta(p)p^{-v}}
 {(1+\eta(p)p^{-v})D_p(s)}.
 \label{jct:bivariate-factor}
\end{equation}
The factors at both deleted rational primes are displayed explicitly.
The local computation proves
\begin{equation}
 \mathcal F(s,v)=A(v)P(s)^{-1}E(s,v)
 \label{jct:bivariate-continuation}
\end{equation}
first in the absolutely convergent region.

The mixed term in \(E\) can be extracted as another Hecke factor.
Write \(x=p^{-s}\), \(y=p^{-v}\), \(e=\eta(p)\), and \(t=t_p\).
At a split prime \(e=1\), and at an inert prime \(t=0\), so
\(et=t\) in both cases. Define
\begin{equation}
 E_{2,p}(s,v)=
 \frac{(1-tx+ey)(1-txy+ex^2y^2)}
 {(1+ey)(1-tx+ex^2)}.
 \label{jct:secondary-local}
\end{equation}
Local multiplication gives, initially in the absolute region,
\begin{equation}
 E(s,v)=P(s+v)E_2(s,v),\qquad
 E_2(s,v)=\prod_{p\ne3,269}E_{2,p}(s,v).
 \label{jct:secondary-factor}
\end{equation}
All local factors at the deleted primes remain omitted in every
product. Subtracting the denominator polynomial in
\eqref{jct:secondary-local} from its numerator gives exactly
\[
 \begin{split}
 -txy^2+x^2\{(t^2-1)y+ey^2+y^3-e\}-tx^3y^2.
 \end{split}
\]
Since \(|t|\le2\) and \(|x|,|y|<1\), this is
\(O(|x|^2+|x||y|^2)\).
The inverse roots of \(D_p\) have modulus one. Thus on
\(\Re s\ge\sigma_0>0\), \(\Re v\ge\nu_0>0\), the denominator
has absolute value at least
\[
 (1-p^{-\nu_0})(1-p^{-\sigma_0})^2,
\]
with a positive lower bound independent of the imaginary parts.
Consequently
\[
 E_{2,p}(s,v)=1+O\bigl(p^{-2\Re s}+p^{-\Re(s+2v)}\bigr).
\]
The product \(E_2\) converges normally, and is uniformly bounded
on closed inner regions, when
\begin{equation}
 \Re s>\tfrac12,\qquad\Re v>0,\qquad\Re(s+2v)>1.
 \label{jct:secondary-domain}
\end{equation}
No nonvanishing assertion for \(E_2\) is needed. Because
\(\beta>1/2\) and \(u\ge1/4\), every point with
\(\Re s>\beta\), \(\Re v>u\) lies strictly inside this domain:
in particular \(\Re(s+2v)>\beta+2u>1\).

\paragraph{Vertical bounds without a simple-zero hypothesis.}
Apply Lemma~\ref{jct:reciprocal-growth} to \(P\) in
\(\Re s>\beta\). Its strip growth is part of the stated Hecke
background, and its absolutely convergent reciprocal Euler product
on \(\Re s=2\) gives the required lower bound there.
The zero-free input consequently yields
\[
 |P(s)^{-1}|\ll_\delta(1+|\Im s|)^{C_\delta}
\]
on every fixed bounded strip with \(\Re s\ge\beta+\delta\).
This argument requires neither a zero-count estimate nor any
restriction on zero multiplicities.

The rational-series argument above gives the corresponding
polynomial bound for \(A(v)\) on every closed strip in
\(\Re v>u=\rho/2\), under the stated zeta zero-free input.
The pole of \(\zeta\) becomes a zero of its reciprocal and causes
no singularity. The additional factor \(P(s+v)\) in
\eqref{jct:secondary-factor} is entire and polynomially bounded
on the fixed strips in question. Its zeros are harmless because
it is not inverted.
Equations~\eqref{jct:bivariate-continuation} and
\eqref{jct:secondary-factor} therefore continue \(\mathcal F\)
holomorphically, with polynomial vertical growth, to
\begin{equation}
 \Re s>\beta,\qquad \Re v>u.
 \label{jct:tube}
\end{equation}
The finite corrections in \eqref{jct:bivariate-factor} have no
zeros in their denominators there.

\paragraph{The two norm caps and the content cutoff.}
Put \(q=U^2/Y\ge1\) and
\[
 \Psi_q(t_1,t_2)=b(t_1)b(t_2)V(qt_1t_2),\qquad
 K(t)=1-\chi(t).
\]
The support of \(\Psi_q\) lies in
\[
 0<t_i\le2,\qquad q^{-1}<t_1t_2<2q^{-1}.
\]
It is compact in the positive quadrant for each fixed \(q\),
since \(t_i>1/(2q)\).  Define
\[
 \widehat\Psi_q(s_1,s_2)=\iint
 \Psi_q(t_1,t_2)t_1^{s_1}t_2^{s_2}
 \frac{dt_1}{t_1}\frac{dt_2}{t_2}.
\]
For fixed \(\sigma>0\), logarithmic coordinates show that the
support has area \(O(\log(2+q))\) and that
\((t_1t_2)^\sigma\ll q^{-\sigma}\).
Each fixed logarithmic derivative of \(\Psi_q\) is bounded
uniformly in \(q\): differentiating \(V(qt_1t_2)\) acts on
its argument.  Repeated integration by parts against the full
Mellin parameters consequently proves, for every integer \(J\ge0\),
\begin{equation}
 |\widehat\Psi_q(\sigma+it_1,\sigma+it_2)|
 \ll_{\sigma,J}q^{-\sigma}\log(2+q)
 (1+|t_1|)^{-J}(1+|t_2|)^{-J}.
 \label{jct:norm-caps}
\end{equation}
For real parts \(\sigma_i\in[\sigma,2]\), the additional
factors \(t_i^{\sigma_i-\sigma}\) are bounded by fixed constants,
so the same estimate controls the intervening contours.
All constants use only finitely many of the stated weight seminorms.

For \(\Re z>0\), define the Mellin transform with the opposite sign
\begin{equation}
 \widetilde K(z)=\int_0^\infty K(t)t^{-z}\frac{dt}{t}
 =\frac1z\int_0^\infty (tK'(t))t^{-z}\frac{dt}{t}.
 \label{jct:cutoff-mellin}
\end{equation}
The boundary terms vanish because \(K=0\) through one and
\(K=1\) from two onward.  Its derivative is smooth and supported
in \([1,2]\).  Hence \(\widetilde K\) is holomorphic for
\(\Re z>0\), and decreases faster than every inverse power of
\(|\Im z|\) on each fixed positive real line.  Mellin inversion gives
\(K(t)=(2\pi i)^{-1}\int_{(\alpha)}\widetilde K(z)t^z\,dz\)
for every fixed \(\alpha>0\).

\paragraph{Successive contour shifts.}
Choose a sufficiently small \(\delta>0\), and put
\[
 \sigma=\beta+\delta,\qquad \nu=u+2\delta,\qquad
 \alpha=2\sigma-\nu=a_\beta>0.
\]
Take \(\delta\) small enough that \(\sigma<2\).
Use the \(z\)-line of real part \(\alpha\) from the outset.
On the initial two norm lines \(\Re s_i=2\), the two values
\(v_i=2s_i-z\) have real part \(4-\alpha>1\).
Absolute convergence, three Mellin inversions and Fubini give
\begin{equation}
\begin{split}
 J_V(Y,U,R)=\frac1{(2\pi i)^3}
 \int_{(\alpha)}\int_{(2)}\int_{(2)}&
 \widetilde K(z)R^{-z}\widehat\Psi_q(s_1,s_2)
 U^{s_1+s_2}\\
 &\cdot\mathcal F(s_1,2s_1-z)
 \mathcal F(s_2,2s_2-z)\,ds_2\,ds_1\,dz.
\end{split}
 \label{jct:triple-inversion}
\end{equation}
Move first the \(s_1\)-line, and then the \(s_2\)-line, to real
part \(\sigma\), keeping the \(z\)-line fixed.  Throughout either
shift, with \(\sigma\le\Re s_i\le2\),
\[
 \Re s_i\ge\sigma>\beta,\qquad
 \Re(2s_i-z)\ge\nu>u.
\]
Thus both factors stay inside \eqref{jct:tube}; no pole is crossed.
Their polynomial bounds depend only polynomially on
\(1+|\Im s_1|+|\Im s_2|+|\Im z|\).
Choose the decay orders in \eqref{jct:norm-caps} and
\eqref{jct:cutoff-mellin} larger than these fixed powers plus four.
The resulting integrable majorant justifies the successive
horizontal limits and all interchanges, including integration in
\(z\).  This is not a simultaneous contour deformation assumption.

On the final lines the scale factors are exactly
\[
 U^{2\sigma}q^{-\sigma}R^{-\alpha}
 =Y^\sigma R^{-\alpha}.
\]
The integral is consequently
\[
 J_V(Y,U,R)\ll Y^{\beta+\delta}
 R^{-a_\beta}\log(2+U^2/Y).
\]
Since \(R\ge1\), take \(\delta\le\varepsilon\) to obtain the
claimed bound.  A nonzero term has content product greater than
\(R\) and norm product at least the square of that content product,
while the shell has norm product less than \(2Y\).
This proves the empty case and completes the general tail theorem,
including its bounded-family uniformity.
